% ID: SEL-D0078-OB-E03
% ESTADO: cotejada. Segunda pasada visual y borrador compilado revisados.
% Fuente: Selectividad/sel_2008_matematicas_aplicadas/Mat CCSS 6 - Examen y criterios de correción.pdf
% SHA-256: 256ad5f6dc324a5bd081279a0455d3ea64e9feb86c8be52341faa83024bc3b5b
% Página: 2; opción B; ejercicio original 3.

\noindent\underline{Parte I}

En una industria de calzado se producen botas y sandalias. De cada 12 pares producidos, 7 pares son botas y 5 de sandalias. La probabilidad de que un par de botas sea defectuoso es $0.08$ y de que lo sea un par de sandalias es $0.03$. Se escoge al azar un par y resulta ser ``no defectuoso''.
\begin{enumerate}[label=\alph*)]
\item \textbf{(1 punto)} \textquestiondown{}Cuál es la probabilidad de que se haya escogido un par de botas?
\item \textbf{(1 punto)} \textquestiondown{}Cuál es la probabilidad de que se haya escogido un par de sandalias?
\end{enumerate}
\par\noindent\begin{minipage}{\linewidth}
\noindent\underline{Parte II}

El consumo, en gramos, de un cierto producto sigue una ley Normal con varianza 225 $\mathrm{g}^2$.
\begin{enumerate}[label=\alph*)]
\item \textbf{(1 punto)} A partir de una muestra de tamaño 25 se ha obtenido una media muestral igual a 175 g. Halle un intervalo de confianza, al $90\%$, para la media del consumo.
\item \textbf{(1 punto)} \textquestiondown{}Cuál debe ser el tamaño mínimo de la muestra para que el correspondiente intervalo de confianza, al $95\%$, tenga una amplitud máxima de 5?
\end{enumerate}
\end{minipage}\par
